\section{Related Work}
\subsection{Graph Retrieval-Augmented Generation}
Graph Retrieval-Augmented Generation (GraphRAG) offers a novel paradigm to enhance the reasoning abilities of LLMs 
by explicitly modeling entity relationships through graph structure such as knowledge graphs~\cite{pengsurvey,hansurvey}. %
GraphRAG has been widely adopted across multiple domains. For instance, GraphRAG based on knowledge graph~\cite{pan2024unifying} incorporates techniques like subgraph retrieval and Graph Neural Networks (GNNs)~\cite{he2024grtriever} to enhance the accuracy and explainability of LLMs in complex knowledge reasoning question answering. GraphRAG based on document graphs~\cite{sonawane2014graph} models semantic relationships between documents and internal document structures, effectively supporting long-document summarization and cross-document information integration~\cite{sarthi2024raptor}. Furthermore, GraphRAG methods for specialized domains are increasingly trending. Examples include leveraging social graphs~\cite{jiang2023social,zeng2024large} to analyze user relationships for precision recommendations and constructing medical or disease graphs~\cite{yang2024kg,li2024dalk} that integrate multi-source heterogeneous data to enable personalized treatment plan generation.

However, existing GraphRAG methods still face significant performance bottlenecks when facing complex reasoning problems.
For example, ToG~\cite{tog,tog2} introduces reasoning steps using breadth-first search (BFS) with pruning approach, but it heavily relies on external LLMs and thus lacks flexibility. 
G-Retriever~\cite{he2024grtriever} adopts K-Nearest Neighbors-based retrieval with limited efficacy in multi-hop scenarios. 
Moreover, existing approaches suffer from shallow retrieval depth, since they often depend on predefined fixed paths or hop counts. 
HippoRAG~\cite{jimenez2024hipporag,gutierrez2025hipporag2} 
only supports single-hop retrieval. 
KGP~\cite{wang2024kgp} and LightRAG~\cite{guo2024lightrag} are typically confined to 2-3 hops, exhibiting performance degradation when handling problems that require deep-dependence mining. 
In contrast, our work improves the reasoning ability of GraphRAG for complex problems using process-constrained RL, aiming to 
effectively resolves the issues of rigid retrieval strategies and complex reasoning shortcomings in existing GraphRAG approaches.

\subsection{RL-Enhanced Retrieval Scheme}
Recently, reinforcement learning (RL) has become increasingly important for complex reasoning~\cite{xie2025logic,yang2024dollmreasoning,ouyang2022training}. %
For example, DeepSeek-R1~\cite{guo2025deepseek} shows that large-scale RL can endow LLMs with deep thinking capabilities. 
Some pioneering works have also adopted RL to enhance RAG~\cite{jin2025searchr1,song2025r1searcher}. For example, R1-Searcher~\cite{song2025r1searcher} and Search-R1~\cite{jin2025searchr1} achieve co-training of retrieval and generation, effectively improving the model's ability to solve complex problems. However, a systematic exploration of RL schemes within the GraphRAG domain remains open. Technically speaking, directly extending existing RL methods into GraphRAG faces several challenges. Firstly, the existing methods usually employ a single outcome-oriented reward function, e.g., based solely on the correctness of final answer. However, this design becomes too simple for GraphRAG, 
since a complex problem may involve complex retrieval calls
of the external knowledge graph. 
More critically, existing methods have often encountered reward hacking in RL training, potentially resulting in shallow retrieval and over-thinking in GraphRAG. To tackle these issues, we propose \methodname~with tailored reward functions and training strategy to non-trivially explore the potential of RL in enhancing GraphRAG.


